\scp{\tsc{\ili{Central} Timor} internal diversity}{05-01}

Each of the \tsc{\ili{Central} Timor} “languages” are internally diverse
with multiple dialects and/or varieties.
\ili{Kemak} speakers report about at l\ili{east} a dozen different named
dialects of their language including: \ili{Atsabe},
\ili{Diirbatin}, \ili{Kailaku}, \ili{Kutubaba}, \ili{Lemia},
\ili{Leolima}, \ili{Leosibe}, \ili{Marobo}, \ili{Oho}, and \ili{Saneri}.
Data for about ten of these varieties is available,
with lexical similarity between any two varieties ranging from
about 76{\%} to 88{\%}.\footnote{
		\ili{Kutubaba} \ili{Kemak} is
		the most divergent known variety of \ili{Kemak}.
		\ili{Kutubaba} \ili{Kemak} is spoken in \ili{Kenebibi} village on the north coast
		of the easternmost \ili{Kemak}-speaking area in \frf{05-01}.}

\citet[i]{LSG2006} lists three different varieties of \ili{Tokodede}:
\ili{Bazartete} \ili{Tokodede}, \ili{Liquiçá} \ili{Tokodede}, and \ili{Maubara} \ili{Tokodede}.
Lexical differences are found between all three varieties as
well as phonological differences,
the most significant of which is \ili{Maubara} \ili{Tokodede} /v/ corresponding
to other \ili{Tokodede} /h/.
One example is *babuy > \ili{Maubara} \ili{Tokodede} \it{vavi}, other \ili{Tokodede} \it{hahi} ‘pig’.

\citet{Fogaca2017} shows that varieties of \tsc{Mambae} fall
into three groupings on the basis of lexical similarity: Northwest
\ili{Mambae}, \ili{Central} \ili{Mambae}, and South \ili{Mambae}.
The South \ili{Mambae} grouping is most divergent with lexical similarity
compared to other varieties of \tsc{Mambae} ranging from 63{\%}--70{\%}.\footnote{
		The degree of lexical similarity between Northwest \ili{Mambae}
		and South \ili{Mambae} is similar to the degree of lexical similarity
		between Northwest \ili{Mambae} and \ili{Tokodede}, which is about 61{\%} \citep[35]{Edwards2019}.}
\citeauthor{Fogaca2017}’s figures summarising the degree of lexical similarity
between different varieties of \tsc{Mambae} are given in \trf{05-02}.


\begin{table}
\caption[\tsc{Mambae} lexical similarity (Fogaça 2017: 82)]
{\tsc{Mambae} lexical similarity \citep[82]{Fogaca2017}\su{†}}\label{tab:05-02}
%\begin{adjustbox}{width=\textwidth}
	\begin{threeparttable} %\st{2pt}
	\begin{tabular}{llllllllllll}
\hhline{---------}
\mc{3}{|l}{Bazartete}	&		&		&		&		&		&	\mc{1}{c|}{}	&&				 						\\
\mc{1}{|l}{81}	&	\mc{3}{l}{Hatulia}				&	\mc{5}{r|}{\tsc{Northwest Mambae}}	&	\\
\mc{1}{|l}{78}	&	81	&	\mc{2}{l}{Railaco}	&		&		&		&		&	\mc{1}{c|}{}	&&	\\ \hhline{-----------}
74	&	75	&	76	&	\mc{3}{|l}{Laulara}	&		&	& & &	\mc{1}{c|}{} \\%&		&		&		&	\\
74	&	76	&	76	&	\mc{1}{|l}{85}	&	\mc{3}{l}{Aileu Vila}	&		&		\mc{3}{r|}{\multirow{2}{*}{\tsc{\ili{Central} Mambae}}}	\\
71	&	73	&	70	&	\mc{1}{|l}{81}	&	83	&	\mc{5}{l}{Hatu-Builico}	&			\mc{1}{c|}{}\\
69	&	69	&	70	&	\mc{1}{|l}{82}	&	80	&	82	&	\mc{4}{l}{Liquidoe}	&		\mc{1}{c|}{} \\ \hhline{~~~---------}
63	&	65	&	65	&	68	&	70	&	68	&	66	&	\mc{3}{|l}{Hatu-Udo}	&		&	\mc{1}{c|}{}\\
63	&	64	&	64	&	67	&	68	&	68	&	67	&	\mc{1}{|l}{82}	&	\mc{2}{l}{Letefoho}	&	\mc{2}{r|}{\tsc{South Mambae}}\\
63	&	66	&	63	&	66	&	68	&	67	&	65	&	\mc{1}{|l}{82}	&	80	&	Betano	& & \mc{1}{c|}{}\\ \hhline{~~~~~~~-----}
		\lspbottomrule
	\end{tabular}
		\begin{tablenotes}\tnsize
			\item[†] {Names of different varieties are from subdistricts where they are spoken.
								Letefoho and Betano are \it{sucos} (“villages”) in the subdistrict of Same.}
			%\item[‡]
			%\item[Δ]
			%\item[‖]
			%\item[¶]
		\end{tablenotes}
	\end{threeparttable}
%\end{adjustbox}
\end{table}

Because this chapter is focused on the internal relationships
of the languages of \tsc{\ili{Central} Timor},
it necessarily highlights the similarities between these languages.
To give a sense of the structural,
phonological, and lexical diversity which is encompassed under
the label “\tsc{Mambae}” numerals in five different varieties
of \tsc{Mambae} are given in \trf{05-03}.

Structural diversity can be seen in the way the numerals 6--9
and multiples of 10 are formed.
While both Northwest \ili{Mambae}
and South \ili{Mambae} have complex forms for the numerals 6--9 (see
\srf{02-03-01-04}), they use different constructions.
Northwest \ili{Mambae} uses \it{hoho} (no known independent meaning
--- also used in \ili{Tokodede} for numerals 6--9).\footnote{
		\ili{Maubara} \ili{Tokodede} uses \it{vovo} for numerals 6--9 pointing to earlier **bobo.}
while South \ili{Mambae} uses \it{liim nai} ‘five \it{nai}’.
\ili{Central} \ili{Mambae} differs again with mono-morphemic forms which
are ultimately from PMP.
For the multiples of ten, three different constructions are found:
Northwest \ili{Mambae} \it{guul} X,
\ili{Central} \ili{Mambae} X \it{nulo} (or metathesised \it{nuul}),
and South \ili{Mambae} \it{saguul haet} X ‘ten times X’ or just \it{haet} ‘times’.
The numerals in \trf{05-03} also give a sample of the phonological
diversity within the \tsc{Mambae} cluster.
Northwest \ili{Mambae} /p/ corresponds to /f/ in \ili{Central} \ili{Mambae}
and either /p/ or /f/ in different varieties of South \ili{Mambae}.
Similarly, \ili{Central} \ili{Mambae} has /k/ where other varieties have /g/.

\begin{table}
	\centering\tcp{Numerals in five varieties of \tsc{Mambae}}{05-03}
	\st{3.9pt}\begin{tabular}{llllll}\lsptoprule
	&	NW. \ili{Mambae}	&	C. \ili{Mambae}	&	C. \ili{Mambae}	&	S. \ili{Mambae}	&	S. \ili{Mambae}\\
	&	\ili{Railaco}	&	Liquidoe	&	Hatu-Builico	&	Hatu-Udo	&	Betano\\ \midrule
1	&	iid	&	iid	&	iid	&	iid	&	iid\\
2	&	ruu	&	ruu	&	ruu	&	rua	&	ruu\\
3	&	teul	&	teul	&	teul	&	teul	&	teul\\
4	&	paat	&	faat	&	faat	&	paat	&	faat\\
5	&	liim	&	liim	&	liim	&	liim	&	liim\\
6	&	hohon iid	&	neen	&	neen	&	liim nai niid	&	liim nai ida\\
7	&	hoho ruu	&	hitu	&	hitu	&	liim nai rua	&	liim nai rua\\
8	&	hoho teul	&	ualu	&	ualu	&	liim nai telu	&	liim nai telu\\
9	&	hoho paat	&	sia	&	sian	&	liim nai pata	&	liim nai fata\\
10	&	saguul	&	sakuul	&	sakuul	&	saguul	&	saguul\\
20	&	guul ruu	&	rua nuul	&	rua nulo	&	saguul haet rua	&	haet rua\\
30	&	guul teul	&	teul nuul	&	teul nulo	&	saguul haet teul	&	haet teul\\
100	&	atus iid	&	atus iid	&	atus iid	&	atus iid	&	atus iid\\
		\lspbottomrule
	\end{tabular}
\end{table}

Further diversity within \tsc{Mambae} can be seen by examining
the data in \citet{Fogaca2017}.
Northwest \ili{Mambae} and \ili{Central} \ili{Mambae} retain a genitive suffix
\it{-n} which has been lost in South \ili{Mambae}; e.g.
Northwest/\ili{Central} \ili{Mambae} \it{ate-n}, South \ili{Mambae} \it{ate} ‘liver’.
Similarly, Northwest \ili{Mambae} and \ili{Central} \ili{Mambae} have a nominal
suffix/enclitic \it{-a} of unclear function; e.g.
Northwest/\ili{Central} \ili{Mambae} \it{ai-a}, South \ili{Mambae} \it{ai} ‘tree’.
\footnote{
		Parallels between Northwest/\ili{Central} \ili{Mambae} \it{-a}
		and the \tsc{Rote-Meto} determiner **-a are striking.
		One function of this determiner in \tsc{Rote-Meto} languages
		is to mark definiteness.}

Another point of diversity is seen in the forms of metathesis.
All varieties of \tsc{Mambae} have productive CV → VC metathesis (e.g.
\it{lelo → leol} ‘sun’),
but different processes of  vowel assimilation operate after metathesis.
Thus, for instance, South \ili{Mambae} from Betano village has lowering
of \it{i → e} after initial \it{a} when metathesis operates,
as seen in *hapuy > *api > \it{aep} ‘fire’.
The varieties of Northwest \ili{Mambae} from \ili{Liquiçá}
and Hatulia have taken this further with complete assimilation
of the initial vowel: *api > \it{eep-a} [ɛpa] ‘fire’.
\ili{Central} \ili{Mambae}, on the other hand,
does not appear to have any vowel assimilation in this particular
form with all known varieties having \it{aif-a} ‘fire’.

Throughout the remainder of the chapter,
we usually cite examples from only one variety each of \ili{Welaun},
\ili{Kemak}, \ili{Tokodede}, Northwest \ili{Mambae}, \ili{Central} \ili{Mambae}, and South \ili{Mambae}.
Unless otherwise stated, \ili{Welaun} data is from the hamlet of Oele{\Q}u
in the village (\it{desa}) of Bauho /bauhoo/.
\ili{Kemak} data is from the \ili{Kailaku} variety,
as spoken in the village of Haikesak/Maumutin.
\ili{Tokodede} data is from the \ili{Liquiçá} /likisaa/ variety, which is
the focus of \citet{LSG2006}.
Northwest \ili{Mambae} data are from the subdistrict of \ili{Railaco},
\ili{Central} \ili{Mambae} from Liquidoe,
and South \ili{Mambae} from the village of Betano in the subdistrict of Same.
